\section{The first initialized training-gate event}
This result concerns the exact binary input dataset and exact archived
initial parameter rows. No numerical trajectory is a proof premise.
Indices in the following statement are zero based.

\begin{proposition}[An isolated transverse first event]
The first training preactivation to vanish is that of neuron $36$ on
sample $3971$. Its time $\tau$ satisfies
\[
 35888/10^{10}<\tau<35891/10^{10}.
\]
The gate changes from inactive to active. Both one-sided normal velocities
are greater than $0.0006064$. All other training gates stay strictly in
their initial signs through the corresponding frozen-mask comparison interval.
The initialized Clarke flow is unique up to this event and has a unique
local continuation across this isolated transverse crossing.
\end{proposition}
\begin{proof}
The interval computation resolves every initial gate strictly. Extend the
initial-mask vector field polynomially and call its unique smooth solution
$y(t)$. Work in the radius-$10^{-6}$ ball about the archived initial state.
The certified speed is at most $0.108613$, so on
$[0,H]$, $H=35891/10^{10}$, a first-exit argument confines $y$ to the ball.
The field Lipschitz and third-trajectory-derivative bounds are constructed
from the full empirical covariance and initial residual. In particular
$\|y'''\|<3.20818$ throughout this interval.

The code evaluates the degree-two Taylor polynomial
\[
 p(t)=y(0)+tF(y(0))+\tfrac12t^2DF(y(0))F(y(0))
\]
with outward interval arithmetic and adds $3.20818t^3/6$ as a norm
remainder. For the candidate preactivation, its upper bound at the lower
time is $-5.3279\cdot10^{-14}$, and its lower bound at the upper time is
$1.2953\cdot10^{-13}$. Its derivative is bounded below by $0.000606439$.
Every other oriented preactivation is positive throughout $[0,H]$:
its initial value plus the negative part of its initial slope times $H$,
minus the uniform quadratic remainder, is greater than
$2.6079\cdot10^{-9}$. Thus the frozen flow has exactly this first event,
and coincides with the initialized flow until it occurs.

At the event the only field jump is the input-coordinate vector
$N^{-1}(w_i^Te_x)x_n$. Outward bounds on the same ball show that adding
its normal projection leaves the post-event normal velocity greater than
$0.000606438$. Both fields therefore cross the isolated hyperplane in the
same direction. Smooth uniqueness on either side and strict transversality
exclude dwelling or a choice of distinct continuations at the boundary.
The full numerical inequalities, exact rational time endpoints, and input
hash are in \texttt{FIRST\_EVENT\_CERTIFICATE.json}.
\end{proof}
This validates one nonlinear event, not a sequence of all subsequent events.
It does not establish initialized entry at time $3.4$.
